% NLA Architecture: Training Pipeline (Stage 0–3)
% % NLA Architecture: Training Pipeline (Stage 0–3)
% % NLA Architecture: Training Pipeline (Stage 0–3)
% % NLA Architecture: Training Pipeline (Stage 0–3)
% \input{latex/04-training}

\begin{frame}{NLA 的核心结构}
\begin{itemize}
  \item 目标：把目标模型某层 activation $h_\ell$ 翻译成自然语言解释 $z$。
  \item AV：Activation Verbalizer，$h_\ell \mapsto z$。
  \item AR：Activation Reconstructor，$z \mapsto \hat h_\ell$。
  \item 训练目标：
  \[
    \min_{\phi,\theta}\;
    \mathbb{E}_{h_\ell \sim \mathcal H}
    \mathbb{E}_{z \sim AV_\phi(\cdot \mid h_\ell)}
    \|h_\ell - AR_\theta(z)\|_2^2 .
  \]
\end{itemize}
\end{frame}

\begin{frame}{NLA RL 迭代}

先做一个监督代理任务：
  \begin{itemize}
    \item AV：activation $\to$ Claude 生成的摘要；
    \item AR：摘要 $\to$ 原 activation。
  \end{itemize}
warm-start 后已有可用信号，FVE 通常约 $0.3$--$0.4$。然后进行迭代：

\begin{enumerate}
  \item 采样一批目标 activation $h_\ell$。
  \item AV 对每个 $h_\ell$ 采样一组解释：
  \[
    z_1,\dots,z_G \sim AV_\phi(\cdot \mid h_\ell).
  \]
  \item AR 根据解释重构 activation $\hat h_i = AR_\theta(z_i)$ ，利用重构误差进行 SFT
  \item 用原始重构误差给 AV 奖励，用 GRPO/RL 更新：
  \[
    r(h_\ell,z_i) = -\|h_\ell - AR_\theta(z_i)\|_2^2.
  \]
\end{enumerate}
\end{frame}

\begin{frame}{退化风险}
\begin{itemize}
  \item 风险一：AV/AR 学出人类看不懂、但 AR 能解码的暗号。
  \item 风险二：AV 直接复读输入上下文，而不是解释 activation。
  \item 论文中的缓解方式：
  \begin{itemize}
    \item SFT warm-start 保持自然语言分布；
    \item \textbf{对 AV 加 KL penalty，约束其接近初始化}；
    \item activation normalization 提高训练稳定性；
    \item 文本瓶颈长度限制，降低完整输入反演能力。
  \end{itemize}
\end{itemize}
\end{frame}


\begin{frame}{NLA 的核心结构}
\begin{itemize}
  \item 目标：把目标模型某层 activation $h_\ell$ 翻译成自然语言解释 $z$。
  \item AV：Activation Verbalizer，$h_\ell \mapsto z$。
  \item AR：Activation Reconstructor，$z \mapsto \hat h_\ell$。
  \item 训练目标：
  \[
    \min_{\phi,\theta}\;
    \mathbb{E}_{h_\ell \sim \mathcal H}
    \mathbb{E}_{z \sim AV_\phi(\cdot \mid h_\ell)}
    \|h_\ell - AR_\theta(z)\|_2^2 .
  \]
\end{itemize}
\end{frame}

\begin{frame}{NLA RL 迭代}

先做一个监督代理任务：
  \begin{itemize}
    \item AV：activation $\to$ Claude 生成的摘要；
    \item AR：摘要 $\to$ 原 activation。
  \end{itemize}
warm-start 后已有可用信号，FVE 通常约 $0.3$--$0.4$。然后进行迭代：

\begin{enumerate}
  \item 采样一批目标 activation $h_\ell$。
  \item AV 对每个 $h_\ell$ 采样一组解释：
  \[
    z_1,\dots,z_G \sim AV_\phi(\cdot \mid h_\ell).
  \]
  \item AR 根据解释重构 activation $\hat h_i = AR_\theta(z_i)$ ，利用重构误差进行 SFT
  \item 用原始重构误差给 AV 奖励，用 GRPO/RL 更新：
  \[
    r(h_\ell,z_i) = -\|h_\ell - AR_\theta(z_i)\|_2^2.
  \]
\end{enumerate}
\end{frame}

\begin{frame}{退化风险}
\begin{itemize}
  \item 风险一：AV/AR 学出人类看不懂、但 AR 能解码的暗号。
  \item 风险二：AV 直接复读输入上下文，而不是解释 activation。
  \item 论文中的缓解方式：
  \begin{itemize}
    \item SFT warm-start 保持自然语言分布；
    \item \textbf{对 AV 加 KL penalty，约束其接近初始化}；
    \item activation normalization 提高训练稳定性；
    \item 文本瓶颈长度限制，降低完整输入反演能力。
  \end{itemize}
\end{itemize}
\end{frame}


\begin{frame}{NLA 的核心结构}
\begin{itemize}
  \item 目标：把目标模型某层 activation $h_\ell$ 翻译成自然语言解释 $z$。
  \item AV：Activation Verbalizer，$h_\ell \mapsto z$。
  \item AR：Activation Reconstructor，$z \mapsto \hat h_\ell$。
  \item 训练目标：
  \[
    \min_{\phi,\theta}\;
    \mathbb{E}_{h_\ell \sim \mathcal H}
    \mathbb{E}_{z \sim AV_\phi(\cdot \mid h_\ell)}
    \|h_\ell - AR_\theta(z)\|_2^2 .
  \]
\end{itemize}
\end{frame}

\begin{frame}{NLA RL 迭代}

先做一个监督代理任务：
  \begin{itemize}
    \item AV：activation $\to$ Claude 生成的摘要；
    \item AR：摘要 $\to$ 原 activation。
  \end{itemize}
warm-start 后已有可用信号，FVE 通常约 $0.3$--$0.4$。然后进行迭代：

\begin{enumerate}
  \item 采样一批目标 activation $h_\ell$。
  \item AV 对每个 $h_\ell$ 采样一组解释：
  \[
    z_1,\dots,z_G \sim AV_\phi(\cdot \mid h_\ell).
  \]
  \item AR 根据解释重构 activation $\hat h_i = AR_\theta(z_i)$ ，利用重构误差进行 SFT
  \item 用原始重构误差给 AV 奖励，用 GRPO/RL 更新：
  \[
    r(h_\ell,z_i) = -\|h_\ell - AR_\theta(z_i)\|_2^2.
  \]
\end{enumerate}
\end{frame}

\begin{frame}{退化风险}
\begin{itemize}
  \item 风险一：AV/AR 学出人类看不懂、但 AR 能解码的暗号。
  \item 风险二：AV 直接复读输入上下文，而不是解释 activation。
  \item 论文中的缓解方式：
  \begin{itemize}
    \item SFT warm-start 保持自然语言分布；
    \item \textbf{对 AV 加 KL penalty，约束其接近初始化}；
    \item activation normalization 提高训练稳定性；
    \item 文本瓶颈长度限制，降低完整输入反演能力。
  \end{itemize}
\end{itemize}
\end{frame}


\begin{frame}{NLA 的核心结构}
\begin{itemize}
  \item 目标：把目标模型某层 activation $h_\ell$ 翻译成自然语言解释 $z$。
  \item AV：Activation Verbalizer，$h_\ell \mapsto z$。
  \item AR：Activation Reconstructor，$z \mapsto \hat h_\ell$。
  \item 训练目标：
  \[
    \min_{\phi,\theta}\;
    \mathbb{E}_{h_\ell \sim \mathcal H}
    \mathbb{E}_{z \sim AV_\phi(\cdot \mid h_\ell)}
    \|h_\ell - AR_\theta(z)\|_2^2 .
  \]
\end{itemize}
\end{frame}

\begin{frame}{NLA RL 迭代}

先做一个监督代理任务：
  \begin{itemize}
    \item AV：activation $\to$ Claude 生成的摘要；
    \item AR：摘要 $\to$ 原 activation。
  \end{itemize}
warm-start 后已有可用信号，FVE 通常约 $0.3$--$0.4$。然后进行迭代：

\begin{enumerate}
  \item 采样一批目标 activation $h_\ell$。
  \item AV 对每个 $h_\ell$ 采样一组解释：
  \[
    z_1,\dots,z_G \sim AV_\phi(\cdot \mid h_\ell).
  \]
  \item AR 根据解释重构 activation $\hat h_i = AR_\theta(z_i)$ ，利用重构误差进行 SFT
  \item 用原始重构误差给 AV 奖励，用 GRPO/RL 更新：
  \[
    r(h_\ell,z_i) = -\|h_\ell - AR_\theta(z_i)\|_2^2.
  \]
\end{enumerate}
\end{frame}

\begin{frame}{退化风险}
\begin{itemize}
  \item 风险一：AV/AR 学出人类看不懂、但 AR 能解码的暗号。
  \item 风险二：AV 直接复读输入上下文，而不是解释 activation。
  \item 论文中的缓解方式：
  \begin{itemize}
    \item SFT warm-start 保持自然语言分布；
    \item \textbf{对 AV 加 KL penalty，约束其接近初始化}；
    \item activation normalization 提高训练稳定性；
    \item 文本瓶颈长度限制，降低完整输入反演能力。
  \end{itemize}
\end{itemize}
\end{frame}
